% from ega4-18-fr.tex
\oldpage[IV-4]{157}

de ces deux propositions sont donc encore valables lorsqu'on suppose l'anneau local $A$ noethérien, hensélien et universellement japonais.
\end{remark}

\begin{corollary}[18.9.11]
\label{IV.18.9.11}
Soient $X$, $Y$ deux préschémas localement noethériens, $Y$ étant supposé géométriquement unibranche et $X$ connexe. Soit $f:X\to Y$ un morphisme réduit (\emph{i.e.} (6.8.1) plat et à fibres géométriquement réduites). Alors, pour toute partie fermée $T$ de $X$ telle que $f(T)$ soit rare dans $Y$, $X-T$ est connexe.
\end{corollary}

Comme $f$ est plat et que $f(T)$, étant rare, ne peut contenir de point maximal de $Y$, $T$ ne peut contenir de point maximal de $X$ (2.3.4), donc est rare dans $X$. Utilisant (15.5.6.1), il suffit de montrer que, pour tout $x\in T$, $\Spec(\sh{O}_{X,x})-\{x\}$ est connexe. Posons $y=f(x)$, qui n'est pas un point maximal de $Y$. Comme l'anneau $\sh{O}_{Y,y}$ est géométriquement unibranche, on peut appliquer (18.9.7) au morphisme
\[
  f_1:\Spec(\sh{O}_{X,x})\longrightarrow\Spec(\sh{O}_{Y,y}),
\]
qui vérifie les conditions (i) et (ii) de ce théorème, et à la partie fermée $\{x\}$ de $\Spec(\sh{O}_{X,x})$, puisque $f_1(\{x\})=\{y\}$ est rare dans $\Spec(\sh{O}_{Y,y})$.

\subsection{Préschémas étales sur un préschéma géométriquement unibranche ou normal}
\EGAReaderTarget{ega4.c04.subsec.18-10}{subsection:IV.18.10}

\begin{theorem}[18.10.1]
\label{IV.18.10.1}
Soient $f:X\to Y$ un morphisme localement de présentation finie, $x$ un point de $X$, $y=f(x)$. Supposons que $Y$ soit intègre et géométriquement unibranche au point $y$. Alors, pour que $f$ soit étale au point $x$, il faut et il suffit que les deux conditions suivantes soient vérifiées:
\begin{enumerate}
  \item[(i)] $f$ est non ramifié au point $x$;
  \item[(ii)] l'homomorphisme $\sh{O}_{Y,y}\to\sh{O}_{X,x}$ est injectif.
\end{enumerate}
De plus, s'il en est ainsi, $X$ est intègre et géométriquement unibranche au point $x$.

Enfin, si $\sh{O}_{Y,y}$ est noethérien, on peut, dans le critère précédent, remplacer la condition (ii) par la condition

\emph{(ii bis)} On a $\dim\sh{O}_{Y,y}\leq\dim\sh{O}_{X,x}$.
\end{theorem}

Le fait que, si $f$ est étale, $X$ soit intègre et géométriquement unibranche au point $x$ est un cas particulier de (17.5.7). Il est clair d'autre part que les conditions (i) et (ii) sont nécessaires pour que $f$ soit étale au point $x$, puisque $\sh{O}_{X,x}$ est alors un $\sh{O}_{Y,y}$-module fidèlement plat (17.6.1). Prouvons que ces conditions sont suffisantes.

Posons $A=\sh{O}_{Y,y}$, $A'={}^{hs}\!A$ (18.8.7), $Y'=\Spec(A')$, $X'=X\times_Y Y'$, $f'=f_{(Y')}:X'\to Y'$. Soient $y'$ le point fermé de $Y'$, $x'$ un point de $X'$ au-dessus de $x$ et de $y'$; puisque le morphisme $g:Y'\to Y$ est plat (18.8.8), tout revient à voir que $f'$ est étale au point $x'$ (17.7.1, (iii)). L'hypothèse (ii) entraîne que le morphisme $\Spec(\sh{O}_{X,x})\to\Spec(\sh{O}_{Y,y})$ est dominant (I, 1.2.7); comme $\sh{O}_{Y,y}$ est intègre, $\Spec(\sh{O}_{Y,y})$ a un seul point générique $y_1$, et il y a donc une générisation $x_1$ de $x$ dans $X$ au-dessus de $y_1$. Le morphisme projection $X'\to X$ étant plat, il existe une générisation $x'_1$ de $x'$ au-dessus de $x_1$ (2.3.4); posons $y'_1=f'(x'_1)$; on a par suite $g(y'_1)=y_1$. Or, l'hypothèse sur $y$ entraîne que $A'$ est intègre (18.8.15), donc $Y'$ a un seul point générique $\eta$, nécessairement au-dessus de $y_1$ puisque $g$

